\section{Algebra}

I dette afsnit vil vi se på, hvordan man kan løse ligninger. At løse en ligning betyder, at man finder værdien af den ukendte variabel, ofte betegnet med $x$.

En ligning er et matematisk udsagn om, at to udtryk er lige store. Et simpelt eksempel er

\[
2 = 2
\]

Denne ligning er sand, fordi begge sider har samme værdi.

Når man arbejder med ligninger, må man udføre de samme operationer på begge sider af ligningen uden at ændre dens sandhedsværdi. For eksempel kan vi gange begge sider med 3:

\[
3\cdot2 = 3\cdot2
\]

hvilket giver

\[
6 = 6
\]

Ligningen er stadig sand, fordi vi har udført den samme operation på begge sider.

Det er også vigtigt at bemærke, at to forskellige udtryk kan repræsentere den samme værdi. For eksempel gælder det, at

\[
3 = \frac{6}{2}
\]

Selvom disse udtryk ser forskellige ud, er de matematisk set ens. Derfor kan vi i en ligning erstatte det ene udtryk med det andet uden at ændre ligningens sandhedsværdi.

Hvis vi for eksempel ganger begge sider af ligningen med 2, kan vi skrive

\[
3\cdot2 = \frac{6}{2}\cdot2
\]

Ved at forenkle højresiden får vi

\[
6 = \frac{6\cdot2}{2} = \frac{12}{2}
\]

Dermed har vi

\[
6 = \frac{12}{2}
\]

Selvom ligningen nu ser anderledes ud end den oprindelige, er den stadig sand, fordi begge sider stadig har samme værdi.

Dette eksempel viser, at man kan omforme ligninger på forskellige måder, så længe man bevarer ligheden mellem de to sider. Denne idé er central i algebra og bruges, når man ønsker at isolere en ukendt variabel i en ligning.

Vi ser her kort på nogle regler for brøkregning.

Når man simplificerer brøker, kan følgende regler være nyttige:

\[
\frac{\frac{a}{b}}{c} = \frac{a}{bc}
\]

For eksempel:

\[
\frac{\frac{1}{2}}{4} = \frac{1}{2\cdot4} = \frac{1}{8}
\]

På samme måde gælder:

\[
\frac{a}{\frac{b}{c}} = \frac{ac}{b}
\]

For eksempel:

\[
\frac{1}{\frac{2}{4}} = \frac{1\cdot4}{2} = \frac{4}{2} = 2
\]

Bemærk, at de to ovenstående regler blot omskriver en brøk til et simplere udtryk. De ændrer ikke på brøkens værdi.

Derfor behøver man ikke ændre på begge sider af en ligning, når man blot simplificerer en brøk, da værdien af brøken ikke ændres. Husk dog, at hvis man udfører en operation, der ændrer værdien på den ene side af en ligning, skal man udføre den samme operation på den anden side for at bevare ligheden.

Betragt for eksempel ligningen

\[
\frac{\frac{16}{8}}{2} = 1
\]

Her kan vi bruge reglen

\[
\frac{\frac{a}{b}}{c} = \frac{a}{bc}
\]

hvor \(a = 16\), \(b = 8\) og \(c = 2\). Vi kan derfor simplificere brøken:

\[
\frac{\frac{16}{8}}{2} = 1
\Rightarrow
\frac{16}{8\cdot2} = 1
\Rightarrow
\frac{16}{16} = 1
\]

Bemærk, at \(\frac{16}{16} = 1\). Når vi simplificerer en brøk, omskriver vi blot udtrykket til en ækvivalent form. Derfor ændrer vi ikke værdien af siden i ligningen.

Når vi ganger to brøker sammen gælder følgende:
\[
\frac{a}{b} \cdot \frac{c}{d} = \frac{ac}{bd} \Rightarrow \frac{1}{2}\cdot\frac{3}{4} = \frac{1\cdot3}{2\cdot4} \Rightarrow \frac{3}{8}
\]
Når vi adderer brøker, skal de have samme nævner. Derfor må vi ofte forkorte eller forlænge en eller flere af brøkerne, før vi kan addere eller subtrahere dem.

Lad os tage et eksempel:

\[
\frac{1}{2} + \frac{3}{4}
\]

Her kan vi bruge vores viden om, at omskrivning af brøker ikke ændrer på deres værdi. Vi ønsker at omskrive $\frac{1}{2}$ til en brøk med samme nævner som $\frac{3}{4}$, altså nævneren 4.

Dette kan vi gøre ved at gange både tælleren og nævneren i $\frac{1}{2}$ med 2. Når vi ganger med det samme tal i både tæller og nævner, ændres værdien af brøken ikke, og den repræsenterer derfor stadig den samme størrelse:

\[
\frac{1\cdot2}{2\cdot2} + \frac{3}{4}
\]

Dette giver

\[
\frac{2}{4} + \frac{3}{4}
\]

Nu har brøkerne samme nævner, og vi kan derfor bruge reglen

\[
\frac{a}{b} + \frac{c}{b} = \frac{a+c}{b}
\]

Ved at anvende denne regel får vi

\[
\frac{2}{4} + \frac{3}{4} = \frac{2+3}{4} = \frac{5}{4}
\]

Vi har altså vist, at

\[
\frac{1}{2} + \frac{3}{4} = \frac{5}{4}
\]

Præcis den samme fremgangsmåde kan anvendes, når man subtraherer brøker. For at kunne løse almene ligninger, mangler vi blot at se på hvoran forkortelse af udtryk skal behandles.

Når vi ønsker at gøre et udtryk større ved at gange med et tal, skal vi være opmærksomme på, hvordan det påvirker brøken. Lad os starte med ligningen:

\[
2 = 2
\]

Dette kan skrives som brøker:

\[
\frac{2}{1} = \frac{2}{1}.
\]

Hvis vi nu ganger 2 på tælleren alene, ændres værdien af brøken. For at bevare ligheden i ligningen skal vi derfor udføre den samme operation på begge sider:

\[
\frac{2\cdot2}{1} = \frac{2\cdot2}{1} \Rightarrow \frac{4}{1} = \frac{4}{1}.
\]

Dette kan forkortes til

\[
4 = 4.
\]

På samme måde gælder det, når vi ønsker at **forkorte** eller dividere. Hvis vi dividerer med 2, skal vi gøre det på begge sider:

\[
\frac{2}{1} = \frac{2}{1} \quad \Rightarrow \quad \frac{2}{1\cdot2} = \frac{2}{1\cdot2} \Rightarrow \frac{2}{2} = \frac{2}{2}.
\]

Dette kan også forkortes til

\[
1 = 1.
\]

Med dette har vi de essentielle regneregler til at arbejde med ligningsløsning:  
\begin{enumerate}
    \item Vi må gange eller dividere begge sider med det samme tal.  
    \item Vi kan omskrive brøker ved at gange eller dividere tæller og nævner, uden at ændre værdien.  
    \item Disse principper sikrer, at ligheden i en ligning altid bevares.
\end{enumerate}
